% ============================================================
%  ARTIPHISHELL REFERENCES
%  Bibliography entries for the Artiphishell paper
%  (DEF CON 2025 / DARPA AIxCC)
%
%  Entries that already exist in frontmatter/references.tex
%  with different keys are aliased here under the paper's
%  original citation keys.
% ============================================================

% ---- Alias entries (paper key differs from existing key) ----

\bibitem{aixcc}
DARPA, ``AIxCC: AI Cyber Challenge,'' 2025, \url{https://aicyberchallenge.com/}.

\bibitem{semgrep}
G. Bennett, T. Hall, E. Winter, and S. Mayfield, ``Semgrep$^+$: Improving the limited performance of static application security testing (SAST) tools,'' in \textit{Int. Conf. on Evaluation and Assessment in Software Engineering (EASE)}, 2024.

\bibitem{aschermann2020ijon}
C. Aschermann, S. Schumilo, A. Abbasi, and T. Holz, ``IJON: Exploring deep state spaces via fuzzing,'' in \textit{2020 IEEE Symposium on Security and Privacy (SP)}, 2020, pp. 1597--1612.

% ---- New entries (not in frontmatter/references.tex) ----

\bibitem{sarif}
OASIS, ``Static Analysis Results Interchange Format (SARIF) Version 2.1.0,'' 2023, \url{https://docs.oasis-open.org/sarif/sarif/v2.1.0/sarif-v2.1.0.html}.

\bibitem{oss-fuzz}
Google, ``OSS-Fuzz,'' 2016, \url{https://google.github.io/oss-fuzz/}.

\bibitem{libfuzzer}
LLVM, ``libFuzzer -- a library for coverage-guided fuzz testing,'' 2025, \url{https://llvm.org/docs/LibFuzzer.html}.

\bibitem{honggfuzz}
R. Swiecki, ``Honggfuzz: Security Oriented Fuzzer with Feedback-Driven Mutations,'' 2019, \url{https://github.com/google/honggfuzz}.

\bibitem{aflplusplus}
A. Fioraldi, D. Maier, H. Eissfeldt, and M. Heuse, ``AFL++: Combining incremental steps of fuzzing research,'' in \textit{Proceedings of the USENIX Conference on Offensive Technologies (WOOT)}, 2020.

\bibitem{jazzer}
Code Intelligence, ``Jazzer: Coverage-Guided, In-Process Fuzzer for the JVM,'' 2025, \url{https://github.com/CodeIntelligenceTesting/jazzer}.

\bibitem{aixcc:toast}
DARPA, ``The Open Application Security Testing (TOAST),'' 2025, \url{https://github.com/aixcc-finals/toast}.

\bibitem{afl}
M. Zalewski, ``American Fuzzy Lop (AFL),'' 2015, \url{https://lcamtuf.coredump.cx/afl/}.

\bibitem{bohme2016aflfast}
M. Böhme, V. Pham, and A. Roychoudhury, ``Coverage-based greybox fuzzing as Markov chain,'' in \textit{Proceedings of the ACM SIGSAC Conference on Computer and Communications Security (CCS)}, 2016.

\bibitem{lyu2019mopt}
C. Lyu, S. Ji, C. Zhang, Y. Li, W. Lee, Y. Song, and R. Beyah, ``MOPT: Optimized mutation scheduling for fuzzers,'' in \textit{Proceedings of the USENIX Security Symposium}, 2019.

\bibitem{aschermann2019redqueen}
C. Aschermann, S. Schumilo, T. Blazytko, R. Gawlik, and T. Holz, ``REDQUEEN: Fuzzing with input-to-state correspondence,'' in \textit{Proceedings of the Network and Distributed System Security Symposium (NDSS)}, 2019.

\bibitem{meng2021:Bran}
D. Meng, M. Guerriero, A. Machiry, H. Aghakhani, P. Bose, A. Continella, C. Kruegel, and G. Vigna, ``Bran: Reduce vulnerability search space in large open-source repositories by learning bug symptoms,'' in \textit{Proceedings of the ACM Asia Conference on Computer and Communications Security (AsiaCCS)}, Hong Kong, China, June 2021.

\bibitem{Moshtari2013}
S. Moshtari, A. Sami, and M. Azimi, ``Using complexity metrics to improve software security,'' \textit{Computer Fraud \& Security}, vol. 2013, no. 5, pp. 8--17, Elsevier, 2013.

\bibitem{Neuhaus2007}
S. Neuhaus, T. Zimmermann, C. Holler, and A. Zeller, ``Predicting vulnerable software components,'' in \textit{Proceedings of the ACM Conference on Computer and Communications Security (CCS)}, 2007.

\bibitem{codeql}
Code Intelligence, ``CodeQL,'' 2025, \url{https://codeql.github.com/}.

\bibitem{nautilus}
C. Aschermann, T. Frassetto, T. Holz, P. Jauernig, A. Sadeghi, and D. Teuchert, ``NAUTILUS: Fishing for deep bugs with grammars,'' in \textit{Proceedings of the Network and Distributed System Security Symposium (NDSS)}, 2019.

\bibitem{benzene}
Y. Park, H. Lee, J. Jung, H. Koo, and H. K. Kim, ``BENZENE: A practical root cause analysis system with an under-constrained state mutation,'' in \textit{Proceedings of the IEEE Symposium on Security and Privacy (SP)}, 2024.

\bibitem{aurora}
T. Blazytko, M. Schlögel, C. Aschermann, A. Abbasi, J. Frank, S. Wörner, and T. Holz, ``AURORA: Statistical crash analysis for automated root cause explanation,'' in \textit{Proceedings of the USENIX Security Symposium}, 2020.

\bibitem{IJON}
C. Aschermann, S. Schumilo, A. Abbasi, and T. Holz, ``IJON: Exploring deep state spaces via fuzzing,'' in \textit{Proceedings of the IEEE Symposium on Security and Privacy (S\&P)}, San Francisco, CA, USA, 2020.

\bibitem{mei2024arvo}
X. Mei, P. S. Singaria, J. Del Castillo, H. Xi, T. Bao, R. Wang, Y. Shoshitaishvili, A. Doup\'{e}, H. Pearce, B. Dolan-Gavitt, and others, ``Arvo: Atlas of reproducible vulnerabilities for open source software,'' \textit{arXiv preprint arXiv:2408.02153}, 2024.

\bibitem{yang2024qwen2}
A. Yang, B. Yang, B. Hui, B. Zheng, B. Yu, C. Zhou, C. Li, C. Chengyuan, D. Liu, F. Huang, and others, ``Qwen2 technical report,'' \textit{arXiv preprint arXiv:2407.10671}, 2024.

\bibitem{grammarinator}
R. Hodován, Á. Kiss, and T. Gyimóthy, ``Grammarinator: A grammar-based open source fuzzer,'' in \textit{Proceedings of the 9th ACM SIGSOFT International Workshop on Automating Test Case Design, Selection, and Evaluation (A-TEST 2018)}, Lake Buena Vista, FL, USA, pp. 45--48, 2018.
